\documentclass[11pt,a4paper]{article}
\usepackage[utf8]{inputenc}
\usepackage[italian]{babel}
\usepackage{geometry}
\usepackage{hyperref}
\usepackage{listings}
\usepackage{xcolor}
\usepackage{amsmath}
\usepackage{graphicx}
\usepackage{enumitem}
\usepackage{titlesec}
\usepackage{fancyhdr}
\usepackage{tikz}
\usetikzlibrary{shapes.geometric, arrows, positioning, fit, backgrounds}

\geometry{margin=2.5cm}

% Configurazione colori
\definecolor{codegreen}{rgb}{0,0.6,0}
\definecolor{codegray}{rgb}{0.5,0.5,0.5}
\definecolor{codepurple}{rgb}{0.58,0,0.82}
\definecolor{backcolour}{rgb}{0.95,0.95,0.92}
\definecolor{primary}{HTML}{6366F1}
\definecolor{secondary}{HTML}{8B5CF6}
\definecolor{accent}{HTML}{F59E0B}
\definecolor{success}{HTML}{10B981}
\definecolor{warning}{HTML}{F59E0B}
\definecolor{danger}{HTML}{EF4444}

\lstset{
    backgroundcolor=\color{backcolour},
    commentstyle=\color{codegreen},
    keywordstyle=\color{magenta},
    numberstyle=\tiny\color{codegray},
    stringstyle=\color{codepurple},
    basicstyle=\ttfamily\footnotesize,
    breakatwhitespace=false,
    breaklines=true,
    captionpos=b,
    keepspaces=true,
    numbers=left,
    numbersep=5pt,
    showspaces=false,
    showstringspaces=false,
    showtabs=false,
    tabsize=2
}

% Header e footer
\pagestyle{fancy}
\fancyhf{}
\fancyhead[L]{\leftmark}
\fancyhead[R]{WorkJournal - Design Document v0.1}
\fancyfoot[C]{\thepage}

\title{\textbf{WorkJournal} \\ \Large Design Document \\ \vspace{0.5cm} \normalsize Riprogettazione Componente Gestione Lavoro}
\author{ExtraTracker - Documentazione Architetturale}
\date{Versione 0.1 - \today}

\begin{document}

\maketitle

\begin{abstract}
Questo documento descrive la progettazione del componente \textbf{WorkJournal}, 
un sistema unificato per la gestione della vita lavorativa quotidiana che combina 
journaling, tracking attività, gestione progetti e supporto AI per prevenire 
stress e burnout. Il documento evolverà incrementalmente durante la fase di design.
\end{abstract}

\tableofcontents
\newpage

%% ============================================================================
%% SEZIONE 1: VISION E OBIETTIVI
%% ============================================================================
\section{Vision e Obiettivi}

\subsection{Mission Statement}
\begin{quote}
\textit{"Rendere la gestione del lavoro quotidiano semplice, chiara e senza stress. 
L'utente deve poter vedere la propria vita lavorativa in pochi sguardi e sapere 
sempre cosa fare, quando farlo, e sentirsi supportato nel processo."}
\end{quote}

\subsection{Obiettivi Primari}

\begin{enumerate}
    \item \textbf{Semplicità d'uso}: Poche azioni con alta resa
    \item \textbf{Chiarezza visiva}: Dashboard che mostra tutto in pochi sguardi
    \item \textbf{Supporto proattivo}: AI che guida e previene burnout
    \item \textbf{Flessibilità}: Adattabile a diversi stili di lavoro
    \item \textbf{Motivazione}: Gamification integrata per engagement
\end{enumerate}

\subsection{Anti-Goals (Cosa NON fare)}

\begin{enumerate}
    \item NON integrare funzionalità di Studio (componente separato esistente)
    \item NON creare interfacce complesse con troppe opzioni
    \item NON richiedere troppi input manuali all'utente
    \item NON duplicare logica tra componenti
\end{enumerate}

\subsection{Principi di Design}

\begin{description}
    \item[Progressive Disclosure] Mostrare prima l'essenziale, dettagli on-demand
    \item[Minimal Input, Maximum Output] L'utente scrive poco, il sistema deduce molto
    \item[AI-First] L'AI non è un add-on, è parte integrante del flusso
    \item[Glanceable] Tutto comprensibile in 3 secondi
\end{description}

\newpage

%% ============================================================================
%% SEZIONE 2: ARCHITETTURA AD ALTO LIVELLO
%% ============================================================================
% --- INIZIO BLOCCO DA COPIARE ---

% Definizioni colori (Rimuovi se già presenti nel tuo preambolo)
\definecolor{primary}{RGB}{70, 130, 180}
\definecolor{secondary}{RGB}{176, 196, 222}
\definecolor{accent}{RGB}{255, 165, 0}
\definecolor{success}{RGB}{60, 179, 113}
\definecolor{warning}{RGB}{255, 99, 71}

\section{Architettura ad Alto Livello}

\subsection{Component Hierarchy}

\begin{figure}[h]
    \centering
    % 'resizebox' forza il contenuto ad adattarsi alla larghezza del testo.
    % È il principio del Responsive Design applicato a LaTeX.
    \resizebox{\textwidth}{!}{%
    \begin{tikzpicture}[
        node distance=1.2cm,
        box/.style={
            rectangle, 
            draw=gray!60, 
            thick,
            rounded corners, 
            minimum width=3.5cm, 
            minimum height=1cm, 
            align=center, 
            fill=white,
            font=\sffamily
        },
        arrow/.style={->, thick, >=stealth, color=gray!70}
    ]

    % Main Container
    \node[box, fill=primary!20, minimum width=12cm, font=\bfseries] (main) {WorkJournalPage};

    % Level 1 - Main Sections (Posizionamento relativo stabile)
    \node[box, fill=secondary!10, below=1.5cm of main] (journal) {JournalSection};
    \node[box, fill=secondary!10, left=1cm of journal] (dashboard) {DailyDashboard};
    \node[box, fill=secondary!10, right=1cm of journal] (projects) {ProjectsOverview};

    % Level 2 - Dashboard children
    \node[box, fill=accent!20, below=1cm of dashboard, font=\small] (insights) {AIInsights};
    \node[box, fill=accent!20, below=0.5cm of insights, font=\small] (alerts) {ProjectAlerts};
    \node[box, fill=accent!20, below=0.5cm of alerts, font=\small] (summary) {YesterdaySummary};

    % Level 2 - Journal children
    \node[box, fill=success!20, below=1cm of journal, font=\small] (quickinput) {QuickInput};
    \node[box, fill=success!20, below=0.5cm of quickinput, font=\small] (timeline) {DayTimeline};
    \node[box, fill=success!20, below=0.5cm of timeline, font=\small] (tasks) {TaskList};

    % Level 2 - Projects children
    \node[box, fill=warning!20, below=1cm of projects, font=\small] (active) {ActiveProjects};
    \node[box, fill=warning!20, below=0.5cm of active, font=\small] (urgent) {UrgentItems};
    \node[box, fill=warning!20, below=0.5cm of urgent, font=\small] (stalled) {StalledProjects};

    % Arrows (Linee ortogonali - Manhattan Style per pulizia visiva)
    \draw[arrow] (main.south) -- ++(0,-0.8) -| (dashboard.north);
    \draw[arrow] (main.south) -- (journal.north);
    \draw[arrow] (main.south) -- ++(0,-0.8) -| (projects.north);

    \draw[arrow] (dashboard) -- (insights);
    \draw[arrow] (insights) -- (alerts);
    \draw[arrow] (alerts) -- (summary);

    \draw[arrow] (journal) -- (quickinput);
    \draw[arrow] (quickinput) -- (timeline);
    \draw[arrow] (timeline) -- (tasks);

    \draw[arrow] (projects) -- (active);
    \draw[arrow] (active) -- (urgent);
    \draw[arrow] (urgent) -- (stalled);

    \end{tikzpicture}
    }
\end{figure}

\subsection{Descrizione Componenti Principali}

\begin{description}
    \item[\texttt{WorkJournalPage}] Container principale, gestisce routing e layout.
    \item[\texttt{DailyDashboard}] Vista ``at-a-glance'' della giornata.
    \item[\texttt{JournalSection}] Area per input giornaliero e timeline.
    \item[\texttt{ProjectsOverview}] Stato progetti e priorità.
\end{description}

\subsubsection{DailyDashboard - Sottocomponenti}
\begin{description}
    \item[\texttt{AIInsights}] Suggerimenti e predizioni dell'AI.
    \item[\texttt{ProjectAlerts}] Progetti in scadenza o che richiedono attenzione.
    \item[\texttt{YesterdaySummary}] Riepilogo automatico di ieri.
\end{description}

\subsubsection{JournalSection - Sottocomponenti}
\begin{description}
    \item[\texttt{QuickInput}] Input rapido per aggiungere entry.
    \item[\texttt{DayTimeline}] Timeline visuale della giornata.
    \item[\texttt{TaskList}] Task/TODO atomici.
\end{description}

\subsubsection{ProjectsOverview - Sottocomponenti}
\begin{description}
    \item[\texttt{ActiveProjects}] Progetti attivi con progresso.
    \item[\texttt{UrgentItems}] Cose urgenti/in scadenza.
    \item[\texttt{StalledProjects}] Progetti fermi da troppo tempo.
\end{description}

\newpage
% ============================================================================
%% SEZIONE 3: DATA MODELS
%% ============================================================================
\section{Data Models (Schema)}

\subsection{Entità Principali}

\subsubsection{Project (Progetto)}

\begin{lstlisting}[language=Java, caption=Project Schema - Pseudocodice]
Project {
    // Identificazione
    id: ObjectId
    name: String                    // "Script AD Martesana"
    code: String                    // "MART-AD-001" (auto o manuale)
    
    // Classificazione
    type: Enum['CLIENT', 'PERSONAL']
    client?: String                 // Nome cliente (se CLIENT)
    category?: String               // Categoria libera
    tags: String[]                  // Tag per ricerca
    
    // Visual
    icon: String                    // Emoji o icona
    color: String                   // Colore hex
    
    // Tracking
    status: Enum['active', 'paused', 'completed', 'archived']
    priority: Enum['low', 'medium', 'high', 'urgent']
    
    // Stime e Budget (opzionali)
    estimatedHours?: Number
    deadline?: Date
    rate?: Number                   // Tariffa oraria (se CLIENT)
    budget?: Number
    
    // Progress
    progress: Number                // 0-100, auto o manuale
    
    // Multi-tenancy
    user: ObjectId
    
    // Timestamps
    createdAt: Date
    updatedAt: Date
}
\end{lstlisting}

\subsubsection{WorkEntry (Entry Giornaliera)}

\begin{lstlisting}[language=Java, caption=WorkEntry Schema - Pseudocodice]
WorkEntry {
    // Identificazione
    id: ObjectId
    project: ObjectId               // Ref -> Project
    
    // Temporalità
    date: Date                      // YYYY-MM-DD
    
    // Tipo di entry
    entryType: Enum[
        'work_block',               // Blocco lavoro con orari
        'note',                     // Nota/riflessione senza orari
        'task_completed',           // Task completato
        'milestone'                 // Milestone raggiunto
    ]
    
    // Contenuto
    title: String                   // Titolo breve
    description?: String            // Descrizione estesa (markdown?)
    
    // Orari (solo per work_block)
    startTime?: String              // "09:00"
    endTime?: String                // "12:30"
    durationMinutes?: Number        // Calcolato o manuale
    
    // Metadata
    tags: String[]
    mood?: Enum['great', 'good', 'neutral', 'tired', 'stressed']
    
    // Gamification
    xpEarned?: Number               // Punti guadagnati
    
    // Multi-tenancy
    user: ObjectId
    
    // Timestamps
    createdAt: Date
    updatedAt: Date
}
\end{lstlisting}

\subsubsection{Task (TODO Atomico)}

\begin{lstlisting}[language=Java, caption=Task Schema - Pseudocodice]
Task {
    // Identificazione
    id: ObjectId
    project?: ObjectId              // Opzionale, può essere generico
    
    // Contenuto
    title: String
    description?: String
    
    // Stato
    status: Enum['pending', 'in_progress', 'completed', 'cancelled']
    priority: Enum['low', 'medium', 'high', 'urgent']
    
    // Scheduling
    dueDate?: Date
    scheduledFor?: Date             // Quando pianificato
    
    // Completamento
    completedAt?: Date
    
    // Gamification
    xpReward: Number                // XP per completamento
    
    // Multi-tenancy
    user: ObjectId
    
    // Timestamps
    createdAt: Date
    updatedAt: Date
}
\end{lstlisting}

\subsection{Relazioni tra Entità}

\begin{center}
\begin{tikzpicture}[
    entity/.style={rectangle, draw, fill=primary!10, minimum width=2.5cm, minimum height=1cm, align=center},
    relation/.style={->, thick, >=stealth}
]

\node[entity] (user) at (0,0) {User};
\node[entity] (project) at (4,0) {Project};
\node[entity] (entry) at (8,0) {WorkEntry};
\node[entity] (task) at (4,-2.5) {Task};

\draw[relation] (user) -- node[above, font=\small] {1:N} (project);
\draw[relation] (project) -- node[above, font=\small] {1:N} (entry);
\draw[relation] (user) -- node[left, font=\small] {1:N} (task);
\draw[relation, dashed] (project) -- node[right, font=\small] {1:N (opt)} (task);

\end{tikzpicture}
\end{center}

\newpage

%% ============================================================================
%% SEZIONE 4: USER FLOWS
%% ============================================================================
\section{User Flows}

\subsection{Flow 1: Apertura Mattutina}

\textbf{Scenario}: L'utente apre WorkJournal la mattina per iniziare la giornata.

\begin{enumerate}
    \item Utente apre \texttt{/workjournal}
    \item Sistema carica:
    \begin{itemize}
        \item Riepilogo AI di ieri
        \item Progetti con alert (scadenze, stalli)
        \item Task pianificati per oggi
    \end{itemize}
    \item Utente vede dashboard "at-a-glance"
    \item AI mostra messaggio tipo: \textit{"Ieri hai lavorato 5h su MART-AD. Oggi potresti completare la fase 2?"}
    \item Utente decide cosa fare e inizia
\end{enumerate}

% TODO: Aggiungere sequence diagram

\subsection{Flow 2: Logging Lavoro Rapido}

\textbf{Scenario}: L'utente vuole registrare un blocco di lavoro appena completato.

\begin{enumerate}
    \item Utente clicca su QuickInput
    \item Scrive: "2h su MART-AD - completato script login"
    \item Sistema:
    \begin{itemize}
        \item Parsa automaticamente durata (2h)
        \item Riconosce progetto (MART-AD)
        \item Estrae descrizione
    \end{itemize}
    \item Conferma con un click
    \item Entry creata + XP assegnati
\end{enumerate}

% TODO: Aggiungere sequence diagram

\subsection{Flow 3: Revisione Fine Giornata}

\textbf{Scenario}: L'utente vuole vedere cosa ha fatto oggi prima di chiudere.

% TODO: Definire questo flow

\subsection{Flow 4: Gestione Progetto}

\textbf{Scenario}: L'utente vuole creare/modificare un progetto.

% TODO: Definire questo flow

\newpage

%% ============================================================================
%% SEZIONE 5: AI INTEGRATION
%% ============================================================================
\section{AI Integration}

\subsection{Ruoli dell'AI}

\begin{enumerate}
    \item \textbf{Analista}: Analizza pattern di lavoro
    \item \textbf{Predittore}: Stima completamento progetti
    \item \textbf{Coach}: Suggerisce bilanciamento e pause
    \item \textbf{Motivatore}: Messaggi di incoraggiamento
\end{enumerate}

\subsection{Dati per l'AI}

L'AI avrà accesso a:

\begin{itemize}
    \item \texttt{WorkEntry[]} - Storico entries
    \item \texttt{Project[]} - Progetti con metriche
    \item \texttt{Task[]} - Task e completamenti
    \item \texttt{UserActivity[]} - Azioni utente (già esistente)
    \item Pattern temporali (ore lavoro per giorno/settimana)
    \item Mood tracking
\end{itemize}

\subsection{Output dell'AI}

% TODO: Definire tipi di output AI

\begin{lstlisting}[language=Java, caption=AIInsight Schema - Pseudocodice]
AIInsight {
    id: ObjectId
    type: Enum[
        'daily_summary',        // Riepilogo giornata
        'prediction',           // Predizione completamento
        'suggestion',           // Suggerimento azione
        'alert',                // Alert burnout/stress
        'motivation'            // Messaggio motivazionale
    ]
    
    content: String             // Testo del messaggio
    
    // Metadata
    relatedProject?: ObjectId
    confidence?: Number         // 0-1, quanto è sicura l'AI
    priority: Enum['low', 'medium', 'high']
    
    // Stato
    dismissed: Boolean
    actedUpon?: Boolean
    
    // Timestamps
    generatedAt: Date
    expiresAt?: Date
}
\end{lstlisting}

\subsection{Trigger AI}

% TODO: Definire quando l'AI genera insights

\newpage

%% ============================================================================
%% SEZIONE 6: GAMIFICATION
%% ============================================================================
\section{Gamification Integration}

\subsection{Punti XP per Azione}

\begin{center}
\begin{tabular}{|l|c|l|}
\hline
\textbf{Azione} & \textbf{XP} & \textbf{Note} \\
\hline
Completare task & 10-50 & Basato su priorità \\
Logging work block & 5 & Base \\
Logging con orari precisi & +5 & Bonus \\
Completare giornata & 20 & Se >= 3 entries \\
Streak giornaliera & 10 * days & Moltiplicatore \\
Raggiungere milestone & 100 & Per progetto \\
\hline
\end{tabular}
\end{center}

\subsection{Achievements Possibili}

% TODO: Definire achievements specifici per WorkJournal

\newpage

%% ============================================================================
%% SEZIONE 7: UI/UX WIREFRAMES
%% ============================================================================
\section{UI/UX Wireframes}

\subsection{Layout Principale}

% TODO: Aggiungere wireframe ASCII o descrizione dettagliata

\begin{lstlisting}[caption=Layout WorkJournal - ASCII Wireframe]
+------------------------------------------------------------------+
|  WORKJOURNAL                              [+ Quick Add]  [Today] |
+------------------------------------------------------------------+
|                                                                   |
|  +------------------------+  +----------------------------------+ |
|  |   DAILY DASHBOARD      |  |      JOURNAL / TIMELINE          | |
|  |                        |  |                                  | |
|  |  [AI Insight Box]      |  |  [Quick Input Field]             | |
|  |  "Ieri hai..."         |  |  ________________________________ | |
|  |                        |  |                                  | |
|  |  [Project Alerts]      |  |  TODAY - Mon 12 Jan              | |
|  |  ! MART-AD scade fra   |  |  +-----------------------------+ | |
|  |    3 giorni            |  |  | 09:00-11:30 | MART-AD       | | |
|  |                        |  |  | Script login completato     | | |
|  |  [Stalled Projects]    |  |  +-----------------------------+ | |
|  |  ~ Progetto X          |  |  | 14:00-?     | CLIENT-Y      | | |
|  |    fermo da 5gg        |  |  | In corso...                 | | |
|  |                        |  |  +-----------------------------+ | |
|  +------------------------+  |                                  | |
|                              |  [Tasks for Today]               | |
|  +------------------------+  |  [ ] Chiamare cliente           | |
|  |   PROJECTS OVERVIEW    |  |  [x] Review PR #123             | |
|  |                        |  |  [ ] Documentazione             | |
|  |  [Active Projects]     |  |                                  | |
|  |  - MART-AD    [====  ] |  +----------------------------------+ |
|  |  - CLIENT-Y   [==    ] |                                      |
|  |                        |                                      |
|  +------------------------+                                      |
+------------------------------------------------------------------+
\end{lstlisting}

\subsection{Responsive Design}

% TODO: Definire comportamento mobile

\newpage

%% ============================================================================
%% SEZIONE 8: API ENDPOINTS
%% ============================================================================
\section{API Endpoints}

\subsection{Endpoints Previsti}

\begin{center}
\begin{tabular}{|l|l|l|}
\hline
\textbf{Method} & \textbf{Endpoint} & \textbf{Descrizione} \\
\hline
GET & /api/workjournal/dashboard & Dashboard data \\
GET & /api/workjournal/entries & Lista entries (filtri) \\
POST & /api/workjournal/entries & Crea entry \\
PUT & /api/workjournal/entries/:id & Modifica entry \\
DELETE & /api/workjournal/entries/:id & Elimina entry \\
\hline
GET & /api/workjournal/projects & Lista progetti \\
POST & /api/workjournal/projects & Crea progetto \\
PUT & /api/workjournal/projects/:id & Modifica progetto \\
\hline
GET & /api/workjournal/tasks & Lista tasks \\
POST & /api/workjournal/tasks & Crea task \\
PATCH & /api/workjournal/tasks/:id/status & Cambia status \\
\hline
GET & /api/workjournal/ai/insights & Insights AI \\
POST & /api/workjournal/ai/generate-summary & Genera summary \\
\hline
\end{tabular}
\end{center}

% TODO: Definire request/response per ogni endpoint

\newpage

%% ============================================================================
%% SEZIONE 9: STATE MANAGEMENT
%% ============================================================================
\section{State Management}

\subsection{Context Structure}

\begin{lstlisting}[language=Java, caption=WorkJournalContext - Pseudocodice]
WorkJournalContext {
    // Data
    projects: Project[]
    entries: WorkEntry[]
    tasks: Task[]
    aiInsights: AIInsight[]
    
    // UI State
    selectedDate: Date
    selectedProject?: string
    viewMode: 'dashboard' | 'timeline' | 'projects'
    
    // Loading States
    loading: {
        projects: boolean
        entries: boolean
        tasks: boolean
        ai: boolean
    }
    
    // Actions
    refreshDashboard(): void
    addEntry(data): Promise<void>
    updateEntry(id, data): Promise<void>
    deleteEntry(id): Promise<void>
    
    addProject(data): Promise<void>
    updateProject(id, data): Promise<void>
    
    addTask(data): Promise<void>
    toggleTaskStatus(id): Promise<void>
    
    dismissInsight(id): void
}
\end{lstlisting}

\newpage

%% ============================================================================
%% SEZIONE 10: SEQUENCE DIAGRAMS
%% ============================================================================
\section{Sequence Diagrams}

% TODO: Aggiungere sequence diagrams per i flow principali

\subsection{Sequence: Caricamento Dashboard}

\begin{lstlisting}[caption=Sequence Diagram - Pseudocodice]
User -> WorkJournalPage: open /workjournal
WorkJournalPage -> WorkJournalContext: mount
WorkJournalContext -> API: GET /api/workjournal/dashboard
API -> ProjectService: getActiveProjects()
API -> WorkEntryService: getRecentEntries()
API -> TaskService: getTodayTasks()
API -> AIService: generateDailyInsights()
API -> WorkJournalContext: { projects, entries, tasks, insights }
WorkJournalContext -> WorkJournalPage: setState()
WorkJournalPage -> User: render dashboard
\end{lstlisting}

\subsection{Sequence: Quick Add Entry}

% TODO: Definire

\subsection{Sequence: AI Insight Generation}

% TODO: Definire

\newpage

%% ============================================================================
%% SEZIONE 11: MIGRATION PLAN
%% ============================================================================
\section{Migration Plan}

\subsection{Fase 1: Preparazione}

\begin{itemize}
    \item [ ] Backup dati esistenti
    \item [ ] Creare nuovi modelli (WorkEntry, Task aggiornati)
    \item [ ] Script migrazione dati da WorkLog -> WorkEntry
\end{itemize}

\subsection{Fase 2: Backend}

\begin{itemize}
    \item [ ] Implementare nuovi service
    \item [ ] Implementare nuovi endpoint
    \item [ ] Implementare AI service base
\end{itemize}

\subsection{Fase 3: Frontend}

\begin{itemize}
    \item [ ] Creare WorkJournalContext
    \item [ ] Implementare componenti base
    \item [ ] Implementare dashboard
    \item [ ] Implementare timeline
\end{itemize}

\subsection{Fase 4: AI Integration}

\begin{itemize}
    \item [ ] Definire algoritmi analisi
    \item [ ] Implementare generazione insights
    \item [ ] Testing e tuning
\end{itemize}

\subsection{Fase 5: Cleanup}

\begin{itemize}
    \item [ ] Rimuovere vecchi componenti PROJECTS
    \item [ ] Rimuovere vecchi componenti WORKSPACE
    \item [ ] Aggiornare routing
    \item [ ] Testing finale
\end{itemize}

\newpage

%% ============================================================================
%% SEZIONE 12: OPEN QUESTIONS
%% ============================================================================
\section{Open Questions}

Domande da risolvere durante la progettazione:

\begin{enumerate}
    \item \textbf{Quick Input Parsing}: Quanto "smart" deve essere? 
    \begin{itemize}
        \item Semplice regex?
        \item NLP con AI?
        \item Template predefiniti?
    \end{itemize}
    
    \item \textbf{Timeline vs Calendar}: Meglio vista timeline verticale o calendario?
    
    \item \textbf{Progetti senza scadenza}: Come gestire progetti senza deadline?
    
    \item \textbf{Multi-device sync}: Real-time sync necessario?
    
    \item \textbf{Offline support}: Necessario?
    
    \item \textbf{Export dati}: Report settimanali/mensili?
    
    % Aggiungi altre domande qui
\end{enumerate}

\newpage

%% ============================================================================
%% APPENDICE: CHANGELOG
%% ============================================================================
\appendix
\section{Changelog}

\begin{description}
    \item[v0.1 - \today] Creazione documento iniziale
    \begin{itemize}
        \item Struttura base documento
        \item Vision e obiettivi
        \item Component hierarchy
        \item Data models base
        \item User flows skeleton
    \end{itemize}
\end{description}

\vspace{2cm}

\begin{center}
\textit{Documento in evoluzione - Versione 0.1}
\end{center}

\end{document}